\documentclass[PMO,lsstdraft,authoryear,toc]{lsstdoc}
\input{meta}

% Package imports go here.
\usepackage{longtable}

% Local commands go here.

%If you want glossaries
%\input{aglossary.tex}
%\makeglossaries

\title{Incident Report}

% This can write metadata into the PDF.
% Update keywords and author information as necessary.
\hypersetup{
    pdftitle={Incident Report},
    pdfauthor={DiegoTapia},
    pdfkeywords={}
}

% Optional subtitle
% \setDocSubtitle{A subtitle}

\input{authors}

\setDocRef{ITTN-086}
\setDocUpstreamLocation{\url{https://github.com/lsst-it/ittn-086}}

\date{\vcsDate}

% Optional: name of the document's curator
% \setDocCurator{The Curator of this Document}

\setDocAbstract{%
Incident Report
}

% Change history defined here.
% Order: oldest first.
% Fields: VERSION, DATE, DESCRIPTION, OWNER NAME.
% See LPM-51 for version number policy.
\setDocChangeRecord{%
  \addtohist{1}{YYYY-MM-DD}{Unreleased.}{Tapia}
}


\begin{document}

% Create the title page.
\maketitle
% Frequently for a technote we do not want a title page  uncomment this to remove the title page and changelog.
% use \mkshorttitle to remove the extra pages

% ADD CONTENT HERE
% You can also use the \input command to include several content files.

\section{Metadata}

\begin{itemize}
  \item \textbf{Incident ID:} 2026-07-02-2
  \item \textbf{Title:} Data deletion in production CNPG PostgreSQL cluster
        (yagan)
  \item \textbf{Date:} 2026-07-02 (detected; data loss occurred 2026-07-01)
  \item \textbf{Duration:} 2026-07-01 (data loss) $\rightarrow$ 2026-07-02
        (detected) $\rightarrow$ 2026-07-03 (restored service verified,
        independent backups live) $\rightarrow$ 2026-07-08 (USDF replication
        restored)
  \item \textbf{Major Incident?} Yes
  \item \textbf{Incident Coordinator:} Diego Tapia
\end{itemize}

\section{Summary}

Data on the production CloudNativePG PostgreSQL cluster on yagan
(\texttt{cnpg-cluster}, namespace \texttt{cloudnativepg}) was discovered
missing on 2026-07-02. The loss was caused by an automatic failover to a stale
replica, triggered when all yagan nodes were hard-reset simultaneously during
scheduled OS updates on 2026-07-01.

Recovery was performed by deploying a parallel CNPG instance
(\texttt{cnpg-cluster-cp}, namespace \texttt{cloudnativepg-cp}) on the same
cluster and restoring the last scheduled backup (2026-07-01 9:00 AM CLT). The
restore completed in 3h 49m with no data-load errors; the restored copy was
verified at row level against production, scaled to a 3-instance HA topology,
and equipped with its own backup pipeline to a dedicated bucket.

\section{Timeline}

\begin{longtable}{p{0.17\textwidth}p{0.76\textwidth}}
\hline
\textbf{Time (CLT)} & \textbf{Event} \\
\hline
\endhead
2026-07-01 9:00 AM & Last scheduled backup dumps production cluster (30 GiB,
  last known-good snapshot) \\
2026-07-01 12:08 PM & All 20 yagan nodes hard-reset simultaneously for OS
  updates via Foreman (\texttt{ipmitool chassis power reset | reboot -h now},
  Evidence D) --- no drain, no switchover. \\
2026-07-01 12:18 PM & Unplanned failover: \texttt{cnpg-cluster-12} promoted to
  primary (16:18:41 UTC, timeline 11 $\rightarrow$ 12; Evidence A), ten
  minutes after the reboot job started. The data loss occurs at this moment. \\
2026-07-02 8:00 AM & Incident detected --- user reports problems with Rapid
  Analysis with the error
  \texttt{lsst.daf.butler.\_exceptions.MissingDatasetTypeError: "Dataset type
  'intrinsicZernikes' does not exist."} Other reports around the same time
  frame of data missing from various other databases. \\
2026-07-02 10:31 AM & Investigation started: inconsistencies found in the CNPG
  cluster via \texttt{kubectl cnpg status cnpg-cluster -n cloudnativepg} ---
  all replicas in file-based standby and behind the primary instance
  (Evidence A). \\
2026-07-02 11:18 AM & Replication normalized on production: stale file-based
  replicas deleted and re-cloned from the primary (548 GB join). By 2:20 PM
  both replicas were streaming, \texttt{Standby (async)}, lag 00:00:00
  (Evidence B). This restored the cluster's health but not its data --- the
  re-cloned replicas were copies of the stale promoted primary, confirming
  recovery required a backup restore, not replication repair. \\
2026-07-02 8:58 PM & Mitigation: parallel-instance recovery chosen (production
  primary left untouched for forensics/rollback); Fleet bundle
  \texttt{cnpg-cluster-new} created and validated on ruka (dev) first. \\
2026-07-02 10:27 PM & \texttt{cnpg-cluster-cp} deployed to yagan (namespace
  \texttt{cloudnativepg-cp}, LB 139.229.180.108, pooler .109). \\
2026-07-02 11:21 PM & Restore started: 2026-07-01 backup dump streamed into
  the new CNPG cluster (\texttt{tar -xOf <backup.tar> | psql}). \\
2026-07-03 3:10 AM & Restore completed after 3h 49m (Evidence C). Log review:
  no data-load errors; only 8 benign ``already exists'' collisions with
  CNPG-provisioned roles/objects. \\
2026-07-03 3:13 PM & Restore verified: missing data confirmed present by the
  system owners. Deployment scaled to 3 instances (primary + 2 streaming
  replicas); DNS record updated to point to the new LB IP. \\
2026-07-03 6:30 PM & Incident resolved: ANALYZE was run on the data and users
  confirmed they could run test images with minimal performance issues
  (Slack). \\
2026-07-08 12:30 PM & USDF replication restored: Pixel firewall rule updated
  to point to the new LB IP; replication slots recreated and resynced with
  USDF. \\
\hline
\end{longtable}

\section{Impact}

\textbf{Affected systems/services:}

\begin{itemize}
  \item Production PostgreSQL on yagan (\texttt{cnpg-cluster}, namespace
        \texttt{cloudnativepg}) and the applications backed by it: ConsDB,
        Butler, NarrativeLog, and Rapid Analysis (which depends on calibration
        data stored in Butler).
\end{itemize}

\textbf{Data loss:}

\begin{itemize}
  \item Apparent loss at failover: several weeks of data --- the database
        reverted to the promoted replica's freeze point.
  \item Permanent loss after restoring the 2026-07-01 9:00 AM dump:
        transactions committed between that backup and the failover, minus
        data reingested afterward. The DM calibration team reingested enough
        calibration data to be ready to go on sky (weather kept the dome
        closed regardless).
\end{itemize}

\textbf{Users affected:} Operations

\textbf{Visibility:} Internal only

\section{Root Cause}

The data loss was caused by an automatic failover to a stale replica,
triggered when all yagan nodes were hard-reset simultaneously during the
2026-07-01 OS updates. The promoted replica is believed to have fallen out of
streaming replication weeks earlier, likely since a previous maintenance.

What is established: on 2026-07-01 at 12:08 PM, all 20 yagan nodes were
hard-reset simultaneously via Foreman (\texttt{ipmitool chassis power reset},
Evidence D) --- a hardware-level power cycle with no node drain, no controlled
switchover, and no graceful OS or PostgreSQL shutdown. Ten minutes later, at
12:18:41 PM, \texttt{cnpg-cluster-12} was promoted to primary, forking the
cluster onto timeline 12 (Evidence A). The scale of the missing data implies
the promoted replica had stopped receiving Write-Ahead Log (WAL) weeks before
the reboot. What remains unclear is when and why it originally lost streaming
replication.

No status capture exists from before the maintenance --- the earliest evidence
is the \texttt{cnpg status} output from the morning of 2026-07-02 (Evidence
A), taken after the failover had already occurred. That snapshot shows the two
remaining replicas in \texttt{Standby (file based)}. On a cluster with WAL
archiving, that state means a replica is catching up from the archive and
typically trails the primary by minutes to hours. Currently our cluster has no
WAL archiving, so a replica that loses streaming receives no WAL at all and
freezes at its last replayed position.

Promotion of a stale replica causes data loss through the following sequence
(standard PostgreSQL/CNPG behavior):

\begin{itemize}
  \item The promoted replica becomes the new primary on a new timeline, forked
        at its last replayed position (LSN). All transactions committed on the
        old primary after that fork point cease to be part of the cluster's
        history.
  \item The remaining replicas re-align to the new primary via
        \texttt{pg\_rewind} or re-clone, discarding any WAL they held beyond
        the fork point. At this step, data that existed on multiple nodes is
        physically removed.
  \item When the former primary rejoined the cluster as a replica, it was
        rewound/re-cloned to follow the new (older) timeline. Until this
        moment, the lost transactions still existed on its disk and were
        recoverable; this step physically deleted the only remaining copy. The
        reboot did not destroy the data --- the automated re-attachment of the
        old primary did.
\end{itemize}

\subsection{Systemic root cause}

The failover-to-stale-replica sequence above is the technical cause, but the
incident was made possible by known, preventable gaps: no replication-state
alerting, no WAL archiving or PITR, no restore runbook, no data-loss
detection, and maintenance tooling that hard-resets database nodes. These gaps
share a single origin: Rubin's DevOps team is severely understaffed relative
to the hundreds of observatory components it operates, and with current
staffing, defensive engineering work of this kind is consistently deferred in
favor of keeping daily operations running. The previous occurrence of this
same failure mode went under-investigated for the same reason, and the same
constraint shaped the response --- a single engineer designed and executed the
entire recovery under pressure, with no second responder available (see
Scorecard).

\section{Resolution}

\begin{enumerate}
  \item Production primary instance left untouched (no in-place restore risk);
        the stale replicas had been deleted and re-cloned earlier during
        investigation to restore streaming replication (Evidence B).
  \item New Fleet bundle \texttt{cnpg-cluster-new} (k8s-cookbook): parallel
        CNPG cluster \texttt{cnpg-cluster-cp} in isolated namespace
        \texttt{cloudnativepg-cp} with dedicated MetalLB IPs
        (139.229.180.108/.109); validated end to end on the ruka dev cluster
        before deploying to yagan.
  \item 2026-07-01 backup restored by streaming the dump:
        \texttt{tar -xOf <backup.tar> | psql -h 139.229.180.108 -U postgres
        -d postgres}. Ran 2026-07-02 11:21 PM $\rightarrow$ 2026-07-03 3:10 AM
        (3h 49m for $\sim$30 GiB). The restore log shows no data-load errors
        --- only 8 benign ``already exists'' collisions with roles/objects
        that CNPG pre-provisions in a fresh cluster (Evidence C).
  \item Restore verified at row level with the system owners of each affected
        database; missing data confirmed present.
  \item Scaled to production-like HA topology: 3 instances (primary + 2
        streaming replicas), pgbouncer pooler, monitoring via PodMonitor.
  \item Dedicated backup pipeline deployed: new S3 bucket and credentials
        (\texttt{cnpg-backups-yagan}), ExternalSecret + CronJob (twice daily),
        verified end to end with a manual run on 2026-07-06.
\end{enumerate}

\section{Lessons Learned}

\textbf{What went well:}

\begin{itemize}
  \item The backup existed, was intact, and restored cleanly --- first full
        restore test of these backups, now proven (Evidence C)
  \item Parallel-instance recovery meant zero additional risk to production,
        and GitOps (Fleet) made it reviewable and reproducible
  \item Dev-first validation (ruka) caught procedural gotchas before touching
        yagan
\end{itemize}

\textbf{What could be improved:}

\begin{itemize}
  \item Detection was manual.
  \item In a data-loss incident, PVCs should be preserved before normalizing
        replication --- the diverged data lives on those disks, and re-cloning
        the replicas during investigation removed potential forensic evidence.
\end{itemize}

\textbf{Gaps identified:}

\begin{itemize}
  \item No point-in-time recovery (WAL archiving not configured)
  \item No data-loss alerting.
  \item No documented restore runbook prior to this incident.
  \item CNPG reported ``Cluster in healthy state'' throughout --- both with
        frozen file-based replicas and after the data loss --- so operator
        status alone cannot detect this condition.
  \item OS-update tooling defaults to hard power resets (Foreman + IPMI) run
        against host groups that include database nodes --- unsafe for
        stateful workloads regardless of replication health.
  \item This happened to a lesser extent in a previous upgrade and a partial
        restore was performed. A more thorough investigation should have been
        done --- effort was lacking however.
\end{itemize}

\section{Preventive actions}

To prevent recurrence of this failure mode, node maintenance involving
database hosts will follow a controlled procedure that guarantees the primary
always steps down cleanly and is never shut down while holding unreplicated
data. Nodes hosting database instances are always maintained one at a time ---
never rebooted simultaneously --- and are restarted with a graceful OS reboot
after the node is drained. Hard resets (\texttt{ipmitool chassis power reset},
Foreman power-cycle jobs) are not used on these nodes except as a last resort
on an already-drained, unresponsive host.

\subsection{Node maintenance procedure for CNPG database clusters}

\textbf{Precondition (gate):} before any node is touched, verify cluster
health with \texttt{kubectl cnpg status}:

\begin{itemize}
  \item Every replica must report \texttt{Standby (async)} (or
        \texttt{Standby (sync)}) --- i.e., attached via \textbf{streaming
        replication} --- with replication lag at or near zero. A replica
        reporting \texttt{Standby (file based)} fails the gate: on clusters
        without WAL archiving (currently all of ours), that state means the
        replica is receiving no WAL at all and is frozen at its last replayed
        position, not ``catching up.''
  \item On clusters where WAL archiving \textbf{is} configured, continuous
        archiving must also be reported healthy (recent successful archive,
        no failing WALs).
\end{itemize}

If any replica fails these checks, maintenance does not proceed until
replication is repaired and lag has returned to $\approx 0$. A lagging or
frozen replica means the cluster has no safe failover target --- an unclean
primary shutdown in this state is exactly the failure mode that caused this
incident.

\begin{enumerate}
  \item \textbf{Replica nodes first.} For each node hosting a replica
        instance:
        \begin{itemize}
          \item Drain the node (\texttt{kubectl cordon} +
                \texttt{kubectl drain}), which shuts the Postgres pod down
                cleanly; the data volume (PVC) remains intact.
          \item Perform the OS update and reboot.
          \item Uncordon the node and wait for the replica to rejoin the
                cluster and fully catch up with the primary (streaming, lag
                $\approx 0$, per \texttt{kubectl cnpg status}) before
                proceeding to the next node.
        \end{itemize}
  \item \textbf{Controlled switchover.} Once all replica nodes are updated,
        promote one of the caught-up replicas (\texttt{kubectl cnpg promote}).
        This is a switchover, not a failover: the old primary confirms a fully
        caught-up successor and hands off cleanly, with no risk of data loss.
        Wait for the switchover to complete and the new primary to report
        healthy.
  \item \textbf{Former primary node last.} Drain the node that hosted the old
        primary (now a replica), perform the OS update and reboot, uncordon,
        and wait for the instance to rejoin and resynchronize with the
        cluster.
\end{enumerate}

At no point in this procedure is a node rebooted while running a database
instance, and the primary role is only ever transferred through a verified,
caught-up successor.

\section{Action Items}

\begin{longtable}{p{0.60\textwidth}p{0.18\textwidth}p{0.10\textwidth}}
\hline
\textbf{Action} & \textbf{Owner} & \textbf{Priority} \\
\hline
\endhead
Alert on any CNPG replica reporting \texttt{Standby (file based)} or
  replication lag above threshold --- all CNPG clusters, all sites &
  [Infrastructure] & High \\
Configure WAL archiving / point-in-time recovery (barman object store) on all
  production CNPG clusters & [Infrastructure] & High \\
Add data-loss alerting to our monitoring stack & [Infrastructure] & High \\
Foreman: separate DB-hosting nodes into their own host group excluded from
  bulk power-cycle jobs; reboots of those nodes follow the CNPG
  node-maintenance runbook & [Infrastructure] & High \\
Publish the node-maintenance procedure above as a runbook and require it in
  the OS-patching procedure for DB-hosting nodes & [Infrastructure] & Medium \\
Write runbook: ``restore a CNPG backup to a parallel instance'' (from IT-6983
  work) & [Infrastructure] & Medium \\
\hline
\end{longtable}

\section{Incident Scorecard}

\begin{longtable}{p{0.12\textwidth}p{0.14\textwidth}p{0.62\textwidth}}
\hline
\textbf{Category} & \textbf{Score} & \textbf{Notes} \\
\hline
\endhead
People & Warning & Single engineer drove the entire recovery; no second
  responder was available for review or relief during a $\sim$20-hour
  incident. \\
Process & Warning & No restore runbook or prior restore test existed; recovery
  procedure had to be designed during the incident and under pressure. \\
Tooling & Warning & Backups + GitOps + operator worked well (OK); but no
  replication-state or data-loss alerting, and maintenance tooling defaults to
  hard resets. \\
\hline
\end{longtable}

\section{References}

\begin{itemize}
  \item \textbf{Related tickets/issues:} IT-6983,
        \url{https://rubinobs.atlassian.net/browse/IT-6983}
  \item \textbf{Communication threads:}
        \href{https://rubin-obs.slack.com/archives/C07RV29CXSR/p1782994199506899}{incident thread}
  \item \textbf{Evidence:} see Appendix (A: cnpg status at detection, B: cnpg
        status after replication normalization, C: restore log, D: Foreman
        reboot job)
\end{itemize}

\appendix

\section{Evidence}

\subsection{A. Production \texttt{cnpg status} at detection (2026-07-02
\textasciitilde10:31 AM CLT)}

Both replicas frozen in file-based standby, $\sim$180--200 GiB of WAL behind
the primary, no WAL archiving --- while the operator reports the cluster
healthy. The promotion time pins the failover to 2026-07-01 16:18:41 UTC
(12:18 PM CLT), ten minutes after the Foreman reboot job (Evidence D) and
$\sim$20 hours before the first user report.

{\footnotesize
\begin{verbatim}
Cluster Summary
Name                 cloudnativepg/cnpg-cluster
System ID:           7533252439990886430
PostgreSQL Image:    docker.io/lsstit/cnpgsphere:16.8
Primary instance:    cnpg-cluster-12
Primary start time:  2026-07-01 16:18:41 +0000 UTC (uptime 21h52m10s)
Status:              Cluster in healthy state
Instances:           3
Ready instances:     3
Size:                549G
Current Write LSN:   159C/3F014960 (Timeline: 12 - WAL File: 0000000C0000159C0000003F)

Continuous Backup status
Not configured

Streaming Replication status
Not available yet

Instances status
Name             Current LSN    Replication role      Status  QoS         Manager Version  Node
----             -----------    ----------------      ------  ---         ---------------  ----
cnpg-cluster-12  159C/3F014960  Primary               OK      Guaranteed  1.27.1           yagan13
cnpg-cluster-13  156F/24000000  Standby (file based)  OK      Guaranteed  1.27.1           yagan14
cnpg-cluster-11  156A/7E93EC00  Standby (file based)  OK      Guaranteed  1.27.1           yagan12
\end{verbatim}
}

\subsection{B. Production \texttt{cnpg status} after replication normalization
(2026-07-02 2:20 PM CLT)}

After deleting the stale replicas, the operator re-cloned them from the
primary (548 GB join per instance). Replication is structurally healthy again
--- but the data is still missing: the replicas are copies of the stale
promoted primary. This output is also the reference ``pass'' state for the
maintenance-gate precondition in Preventive Actions.

{\tiny
\begin{verbatim}
Cluster Summary
Name                     cloudnativepg/cnpg-cluster
System ID:               7533252439990886430
PostgreSQL Image:        docker.io/lsstit/cnpgsphere:16.8
Primary instance:        cnpg-cluster-12
Primary promotion time:  2026-07-01 16:18:41 +0000 UTC (26h1m47s)
Status:                  Cluster in healthy state
Instances:               3
Ready instances:         3
Size:                    549G
Current Write LSN:       159C/7200ED08 (Timeline: 12 - WAL File: 0000000C0000159C00000072)

Continuous Backup not configured

Streaming Replication status
Replication Slots Enabled
Name             Sent LSN       Write LSN      Flush LSN      Replay LSN     Write Lag  Flush Lag  Replay Lag  State      Sync State  Sync Priority  Replication Slot
----             --------       ---------      ---------      ----------     ---------  ---------  ----------  -----      ----------  -------------  ----------------
cnpg-cluster-14  159C/7200ED08  159C/7200ED08  159C/7200ED08  159C/7200ED08  00:00:00   00:00:00   00:00:00    streaming  async       0              active
cnpg-cluster-15  159C/7200ED08  159C/7200ED08  159C/7200ED08  159C/7200ED08  00:00:00   00:00:00   00:00:00    streaming  async       0              active

Instances status
Name             Current LSN    Replication role  Status  QoS         Manager Version  Node
----             -----------    ----------------  ------  ---         ---------------  ----
cnpg-cluster-12  159C/7200ED08  Primary           OK      Guaranteed  1.27.1           yagan13
cnpg-cluster-14  159C/7200ED08  Standby (async)   OK      Guaranteed  1.27.1           yagan14
cnpg-cluster-15  159C/7200ED08  Standby (async)   OK      Guaranteed  1.27.1           yagan12
\end{verbatim}
}

\subsection{C. Restore log (\texttt{restore-yagan.log})}

Full psql output of the dump restore into \texttt{cnpg-cluster-cp}, 2026-07-02
23:21 $\rightarrow$ 2026-07-03 03:10 (3h 49m, $\sim$30 GiB). No data-load
errors; 8 benign ``already exists'' collisions with CNPG-provisioned
roles/objects.

\subsection{D. Foreman job --- yagan node reboots (2026-07-01)}

Job ``\texttt{Run ipmitool chassis power reset | reboot -h now}'', started
2026-07-01 12:08 PM (GMT-4), effective user root, template ``Run Command --
Script Default'', targeting all 20 yagan hosts simultaneously. Result: 16
succeeded, 4 failed (yagan13, yagan17, yagan18, yagan19). The unplanned
primary promotion (Evidence A) occurred at 12:18:41 PM CLT --- ten minutes
after the job started.

\section{Acknowledgements}

This material is based upon work supported in part by the National Science Foundation through Cooperative Agreements AST-1258333 and AST-2241526 and Cooperative Support Agreements AST-1202910 and AST-2211468 managed by the Association of Universities for Research in Astronomy (AURA), and the Department of Energy under Contract No.\ DE-AC02-76SF00515 with the SLAC National Accelerator Laboratory managed by Stanford University.
Additional Rubin Observatory funding comes from private donations, grants to universities, and in-kind support from LSST-DA Institutional Members.

% Include all the relevant bib files.
% https://lsst-texmf.lsst.io/lsstdoc.html#bibliographies
\section{References} \label{sec:bib}
\renewcommand{\refname}{} % Suppress default Bibliography section
\bibliography{local,lsst,lsst-dm,refs_ads,refs,books}

% Make sure lsst-texmf/bin/generateAcronyms.py is in your path
\section{Acronyms} \label{sec:acronyms}
\input{acronyms.tex}
% If you want glossary uncomment below -- comment out the two lines above
%\printglossaries





\end{document}
